\begin{figure} [h]
	
	\centering
	
	\tikzstyle{comp} = [circle, draw, minimum size = 0.5cm, line width=0.7] % comparator style
	\tikzstyle{block} = [rectangle, draw, minimum width = 1.2cm, minimum height = 1.2cm, line width=0.7]  % bloc style
	\tikzstyle{big_block} = [rectangle, draw, minimum width = 2cm, minimum height = 2cm, line width=0.7]  % bloc style
	
	\newcommand{\comp}[3][]{ % comparator command
		\node[comp] (#2) at #3 {};
		\draw (#2) -- + (45:0.25) -- + (45:-0.25);
		\draw (#2) -- + (-45:0.25) -- + (-45:-0.25);
		\ifthenelse{
			\isempty{#1}{}
		}{}
		{
			\ifthenelse{
				\equal{#1}{a}
			}{
				\node[above right] at (#2) {$-$\phantom{p}};
			}{}
			\ifthenelse{
				\equal{#1}{b}
			}{
				\node[below right] at (#2) {$-$\phantom{l}};
			}{}
		}
	}
	\resizebox{\textwidth}{!}{
	\begin{tikzpicture} [node distance = 2cm]
	
	\comp[a]{c1}{(0, 0)};
	
	\node[left of = c1, node distance = 1.5cm] (tau_r) {};
	
	\node[block, right of = c1] (imp_ctrl) {$Cor.$};
	
	%\comp[b]{c2}{($(imp_ctrl) + (2.5,0)$)};
	
	\node[block, above of = imp_ctrl, align = center, node distance = 2cm, xshift=1cm] (transf) {transformation};
	%\node[block, right of = neg, node distance = 2.5cm](jacobian) {$\jacobian{q}{}{t}$};
	
	%\node[block, right of = c2] (kv) {$\vect{K}_v$};
	
	\node[big_block, right of = imp_ctrl, node distance = 2.7cm] (robot) {$Robot$};
	
	\node[right of = robot, node distance = 1.75cm] (force) {};
	\fill ($ (force) + (0, 0.5) $) circle[radius=1.75pt];
	
	\node[big_block, right of = robot, node distance = 3.5cm] (env) {$-Z_{env}$};
	
	\draw[{Latex[length=3mm]}-] (c1.west) --+ (-1.5,0)
		node[align = center, above] {$\vect{\tau}_r = 0$};
	\draw[{Latex[length=3mm]}-] (c1.north) |- (transf.west);
	\draw[-{Latex[length=3mm]}] (c1.east) -- (imp_ctrl.west);
	
	%\draw[-{Latex[length=3mm]}] (imp_ctrl.east) -- (c2.west)
	%node[midway, align = center, above] {$\dot{\vect{q}}_d$};
	%\draw[-{Latex[length=3mm]}] (c2.east) -- (kv.west);
	
	%\draw[-{Latex[length=3mm]}] (kv.east) -- (robot.west)
	%node[midway, above] {$\vect{\tau}_d$};
	\draw[-{Latex[length=3mm]}] (imp_ctrl.east) -- (robot.west)
		node[midway, above] {$\dot{\vect{q}}_d$};
	
	%\draw[{Latex[length=3mm]}-] (neg.east) -- (jacobian.west);
	
	%\draw[-{Latex[length=3mm]}] (robot.south) --+ (0,-0.5)  
	%	node[below] {$\dot{\vect{q}}_m$} -| (c2.south);
	
	\draw[-{Latex[length=3mm]}] ($ (force.center) + (0, 0.5) $) |- (transf.east);
	\draw[{Latex[length=3mm]}-] ($(robot.east) + (0, 0.5) $) -- ($(env.west) + (0, 0.5) $)
		node[midway, align = center, below] {$\vect{f}$};
	\draw[-{Latex[length=3mm]}] ($(robot.east) + (0, -0.5) $) -- ($(env.west) + (0, -0.5) $)
		node[midway, align = center, below] {$\vect{v}$};
	
	\draw[gray, -{Latex[length=3mm]}] (env.south west) -- (env.north east);
	
	\end{tikzpicture}
	}
	\caption{Schéma simplifié du contrôle en transparence avec capteur de force proposé par \cite{lamy_achieving_2009}, avec Cor. un correcteur permettant de générer un signal de commande en vitesse articulaire $\dot{\vect{q}}_d$, et $Z_{env}$ l'impédance de l'environnement.} % Figure caption
		%Bloc diagram of the admittance control with a force sensor, proposed by \cite{Lamy_achieving_2009}}
		%
	\label{lamy_control} % Label 
\end{figure}